

% !TEX root = ./mathcraft.tex
%%%%%%%%%%%%%%%%%%%%%%%%%%%%%%%%%%%%%%%%%%%%%%%%%%%%%%%%%%%
% Obere Titelmakros. Editieren Sie diese Datei nur, wenn
% Sie sich ABSOLUT sicher sind, was Sie da tun!!!
% Version 2.2 - 2.3.2010
%%%%%%%%%%%%%%%%%%%%%%%%%%%%%%%%%%%%%%%%%%%%%%%%%%%%%%%%%%%
\newcommand{\AN}{twoside}
\newcommand{\AUS}{}


%\newcommand{\englisch}{}
%\newcommand{\deutsch}{\usepackage[german]{babel}}

%% Die folgenden auskommentierten Optionen dienen der automatischen
%% Erkennung des Latex-Kompilers und dem Setzen der davon abhängigen
%% Einstellungen. Bei Problem z.B. mit dem Einbinden von verschiedenen
%% Grafiktypen bei Verwendung von PdfLatex oder Latex, einfach die
%% verschiedenen \usepackage(s) ausprobieren. (Mit diesen Einstellungen
%% funktionierte diese Vorlage bei der Verwenundg von latex.exe als
%% Kompiler bei den meisten Studierenden.)

%\newif\ifpdf \ifx\pdfoutput\undefined
%\pdffalse % we are not running pdflatex
%\else
%\pdfoutput=1 % we are running pdflatex
%\pdfcompresslevel=9 % compression level for text and image;
%\pdftrue \fi

\documentclass[11pt,a4paper, \zweiseitig]{article}
\usepackage{ifthen}


%\usepackage[iso]{umlaute}
\usepackage{iftex}
\ifPDFTeX
	\PackageError{imports}{This document must be compiled with XeLaTeX (or LuaLaTeX)}{Run xelatex (see the recipe in .vscode/settings.json or .latexmkrc), not pdflatex.}
\fi
% inputenc/fontenc are for pdfLaTeX only; XeLaTeX reads UTF-8 natively and
% fontspec replaces the OT1/T1 font-encoding machinery.
\usepackage{fontspec}
\defaultfontfeatures{Ligatures=TeX}
\setmainfont{TeX Gyre Pagella} % OpenType clone of Palatino, replaces \usepackage{palatino}
% Monospace font used for \texttt and code listings (minted below). JuliaMono
% is designed for scientific/programming use and has by far the broadest
% Unicode coverage of any monospace font (math operators, arrows, blackboard
% bold, etc.), which is exactly the symbol set used in the pseudocode in this
% document (▶, ×, ℕ, Ω, λ, ∀, ≤, ...). This lets code blocks use those symbols
% as plain UTF-8 characters instead of dropping into \mathescape.
\setmonofont{JuliaMono}[Scale=MatchLowercase]
% Fallback font for symbols/emoji outside the main font (e.g. \emoji{🌳}).
% Symbola covers a very large swath of Unicode, incl. most emoji, in outline
% (monochrome) form, which is what XeLaTeX can actually typeset; the color
% "Noto Color Emoji" font will not render in color under XeLaTeX/LuaLaTeX.
\newfontfamily\emojifont{Symbola}
\newcommand{\emoji}[1]{{\emojifont #1}}
%\usepackage{makeidx} % um ein Index zu erstellen
\usepackage[nottoc]{tocbibind}
\usepackage{fancybox} % fuer die Rahmen
\usepackage{shortvrb}
\usepackage{url}
\usepackage{xcolor}
\usepackage{nicefrac}
% minted (https://ctan.org/pkg/minted) replaces `listings`: it renders code
% via Pygments, so any Unicode symbol supported by the chosen font (see
% \setmonofont above) can be typed directly in the listing, and it produces
% much nicer highlighted/framed output ("beautiful pseudocode") than
% `listings`. It shells out to Pygments (via the `latexminted` helper), so
% compilation needs -shell-escape (already set in .latexmkrc / .vscode).
\usepackage{minted}
\usemintedstyle{bw} % Pygments "bw": bold/italic only, no color or background
\setminted{
	fontsize=\small,
	breaklines=true,
	autogobble=true, % strip common leading whitespace, so indented \begin{minted} blocks aren't shifted right
	linenos=true, % classic look: line numbers in the margin,
	frame=leftline, % separated from the code by a vertical rule
	framesep=6pt,
	numbersep=8pt,
	ignorelexererrors=true, % Pygments flags non-ASCII like ▶ as errors and draws a red box around them
	baselinestretch=1.2,
}
\renewcommand{\theFancyVerbLine}{\scriptsize\arabic{FancyVerbLine}}
\setlength{\marginparwidth}{2cm}
\usepackage{todonotes}
\usepackage{enumitem} % (a), (b), ... labels in definitions
\usepackage[colorlinks,citecolor=black,linkcolor=black,urlcolor=black]{hyperref} %anklickbares Inhaltsverzeichnis

\ifthenelse{\equal{\biber}{true}}{
	% only needed for biber
	\usepackage{csquotes}
	\usepackage[style=numeric,sorting=none,natbib=true,backend=biber,mincitenames=1,maxcitenames=2,maxbibnames=99,uniquelist=false]{biblatex}

	% https://tex.stackexchange.com/a/334703/8850
	\AtEveryBibitem{%
		\clearfield{issn}
		\clearfield{isbn}
		\clearfield{doi}
		\clearfield{location}
		\clearlist{location}
		\clearlist{address}

		\ifentrytype{online}{}{% Remove url except for @online and @thesis
			\ifentrytype{thesis}{}{\clearfield{url}}
		}
	}
}
{}%no else

% Falls es bei \citet ein Komma zwischen Name und Jahr gibt:
% https://tex.stackexchange.com/questions/312539/unwanted-comma-between-author-and-year-using-citet-command
% (thx @ Markus Brenneis)
%\DeclareDelimFormat[cbx@textcite]{nameyeardelim}{\addspace}



\ifthenelse{\equal{\sprache}{deutsch}}{
	\usepackage[ngerman]{babel}
	% Bibtex u.a -> et al.
	\ifthenelse{\equal{\biber}{true}}{
		\DefineBibliographyStrings{ngerman}{
			andothers = {{et\,al\adddot}},
		}
		\newcommand{\references}{Literatur}
	}
	{} % do nothing when not using biber
	\usepackage[autostyle, german=quotes]{csquotes} % Deutsche Anführungszeichen im Literaturverzeichnis (thx @ Markus Brenneis)

}{ \newcommand{\references}{References}}

\usepackage{a4wide} % ganze A4 Weite verwenden



%\ifpdf
%\usepackage[pdftex,xdvi]{graphicx}
%\usepackage{thumbpdf} %thumbs fuer Pdf
%\usepackage[pdfstartview=FitV]{hyperref} %anklickbares Inhaltsverzeichnis
%\else
%\usepackage[dvips,xdvi]{graphicx}
\usepackage{graphicx}
\usepackage{caption}
\usepackage{subcaption}
\usepackage[section]{placeins} % floats stay within their section
\usepackage{array}
\usepackage{booktabs}

%\fi

% Mathematik / kategorientheoretische Notation und kommutative Diagramme
\usepackage{amsmath}
\usepackage{mathtools} % \coloneqq
\usepackage{amssymb}
\usepackage{amsthm}
\usepackage{stmaryrd} % \llbracket, \rrbracket
\usepackage{tikz-cd}
\usepackage{tikz-3dplot}
\usepackage{pgfplots}
\usepackage[super]{nth}
\pgfplotsset{compat=1.18}
\usetikzlibrary{3d, positioning, arrows.meta, patterns, patterns.meta}
% chains of nodes drawn by the forests F(U_m) of the elements U_m of \Omega(n)
\tikzset{forestrow/.style={baseline=-0.5ex, x=0.6cm, >=Stealth},
	onode/.style={draw, rounded corners=3pt, fill=white, inner sep=1pt,
			font=\scriptsize, minimum width=4mm, minimum height=3.5mm},
	oedge/.style={->, thin}}
% \forestchain{0,1,\cdots,n}: a chain of labeled nodes, each pointing to its
% left neighbour; a \cdots entry is drawn without a box
\def\forestchaindots{\cdots}
\newcommand{\forestchain}[1]{\tikz[forestrow]{%
		\foreach \l [count=\i, remember=\i as \j (initially 0)] in {#1} {%
				\ifx\l\forestchaindots \tikzset{chainnode/.style={inner sep=0pt}}%
				\else \tikzset{chainnode/.style={onode}}\fi
				\ifnum\i=1 \node[chainnode] (v\i) at (0, 0) {$\l$};%
				\else \node[chainnode, right=1.8mm of v\j] (v\i) {$\l$}; \draw[oedge] (v\i) -- (v\j);\fi
			}}}
% inline glyphs for the node labels of the graphs in Fig.~\ref{fig:lattices},
% drawn like the nodes there: a cube, a face, an edge and a vertex
\tikzset{cellglyphpic/.style={x={(1.0cm,0cm)}, y={(0.55cm,0.38cm)}, z={(0cm,1.0cm)},
			scale=0.26, line join=round, line cap=round, baseline=0.1ex},
	inlinecellglyph/.style={line width=0.45pt, preaction={fill=white}}}
\newcommand{\cubeglyph}{\tikz[cellglyphpic]{%
		\draw[inlinecellglyph] (0,0,0) -- (1,0,0) -- (1,0,1) -- (0,0,1) -- cycle;
		\draw[inlinecellglyph, pattern={Lines[angle=90, distance=1.1pt, line width=0.3pt]}]
		(1,0,0) -- (1,1,0) -- (1,1,1) -- (1,0,1) -- cycle;
		\draw[inlinecellglyph, pattern={Dots[distance=1.2pt, radius=0.3pt]}]
		(0,0,1) -- (1,0,1) -- (1,1,1) -- (0,1,1) -- cycle;}}
\newcommand{\faceglyph}{\tikz[cellglyphpic]{%
		\draw[inlinecellglyph] (0,0,0) -- (1,0,0) -- (1,0,1) -- (0,0,1) -- cycle;}}
\newcommand{\edgeglyph}{\tikz[cellglyphpic]{%
		\draw[line width=2.6pt, line cap=round] (0,0,0.5) -- (1,0,0.5);
		\draw[line width=1.4pt, line cap=round, white] (0,0,0.5) -- (1,0,0.5);}}
\newcommand{\vertexglyph}{\tikz[baseline=0.1ex]{%
		\draw[inlinecellglyph] (0,0.13cm) circle[radius=2.2pt];}}
\usepackage{mathpartir} % inference rules for the internal logic
\usepackage{aliascnt}
\usepackage[capitalise,nameinlink]{cleveref}

\newtheorem{theorem}{Theorem}[section]
% each environment gets an alias of the theorem counter, so that the shared
% numbering is kept but cleveref sees the environment name, not "theorem"
\newcommand{\newsharedtheorem}[2]{%
	\newaliascnt{#1}{theorem}%
	\newtheorem{#1}[#1]{#2}%
	\aliascntresetthe{#1}}
\newsharedtheorem{definition}{Definition}
\newsharedtheorem{proposition}{Proposition}
\newsharedtheorem{lemma}{Lemma}
\newsharedtheorem{corollary}{Corollary}
\newsharedtheorem{example}{Example}
\newsharedtheorem{remark}{Remark}

\newcommand{\Set}{\mathbf{Set}}
\newcommand{\Top}{\mathbf{Top}}
\newcommand{\functorcat}[2]{[#1,\, #2]}
\newcommand{\toposoftrees}{\mathcal{S}}
\newcommand{\toposofgraphs}{\mathcal{G}}
\newcommand{\restr}[2]{#1|_{#2}}
\newcommand{\Sub}[1]{\mathrm{Sub}_{#1}}
\newcommand{\NN}{\mathbb{N}}
\newcommand{\ZZ}{\mathbb{Z}}
\newcommand{\RR}{\mathbb{R}}
\newcommand{\CC}{\mathbb{C}}
\newcommand{\Stream}{\texttt{Str}}
\newcommand{\List}{\mathrm{List}}
\newcommand{\Blocks}{\mathrm{Blocks}}
\newcommand{\Entities}{\mathrm{Entities}}
\newcommand{\BlockID}{\mathrm{BlockID}}
\newcommand{\EntityID}{\mathrm{EntityID}}
\newcommand{\Hom}{\text{Hom}}
\newcommand{\Nat}{\text{Nat}}
\newcommand{\StreamG}{\Stream}
\newcommand{\StreamedTreeG}{\texttt{StrTree}}
\newcommand{\BinTreeG}{\texttt{Tree}}
\newcommand{\im}{\text{im }}
\newcommand{\treeEmbedding}{F}
\newcommand{\later}{\mathord{\blacktriangleright}}
\newcommand{\before}{\mathord{\blacktriangleleft}}
\newcommand{\laterModal}{\triangleright}
\newcommand{\alwaysModal}{\Box}
\newcommand{\eventuallyModal}{\Diamond}
\newcommand{\liftingMonad}{L}
\newcommand{\laterT}{\triangleright}
\newcommand{\fixLater}{\text{fix}}
% semantic brackets for the interpretation of the internal language
\newcommand{\sem}[1]{[\![#1]\!]}
\newcommand{\id}{\text{id}}
\newcommand{\avg}{\text{avg}}
\newcommand{\true}{\text{true}}
\newcommand{\false}{\text{false}}
\newcommand{\Hask}{\mathbf{Hask}}
\newcommand{\Bool}{\text{Bool}}
\newcommand{\RForest}{\mathbf{RForest}}
% Signature components as in the Elephant (D1.1.1); optional argument is the signature
\newcommand{\SigType}[1][\Sigma]{#1\text{-}\mathrm{Type}}
\newcommand{\SigFun}[1][\Sigma]{#1\text{-}\mathrm{Fun}}
\newcommand{\SigRel}[1][\Sigma]{#1\text{-}\mathrm{Rel}}
\newcommand{\SigTerm}[1][\Sigma]{#1\text{-}\mathrm{Term}}
\newcommand{\SigForm}[1][\Sigma]{#1\text{-}\mathrm{Form}}

\newcommand{\redt}[1] {
	\textcolor{red}{#1}}

\newcommand{\oranget}[1] {
	\textcolor{orange}{#1}}

\newcommand{\purplet}[1] {
	\textcolor{purple}{#1}}

%%%%%%%%%%%%%%%%%%%%%%% Massangaben fuer die Arbeit %%%%%%%%%%%%%%%
\setlength{\textwidth}{15cm}

\setlength{\oddsidemargin}{35mm}
\setlength{\evensidemargin}{25mm}

\addtolength{\oddsidemargin}{-1in}
\addtolength{\evensidemargin}{-1in}

\ifthenelse{\equal{\biber}{true}}{\addbibresource{references.bib}}{}

%\makeindex
